\section{References and retained evidence}
\begin{enumerate}
\item ProveIt, main commit
\texttt{9cca523d9f27ad39d52993f4065cbd7e58f8bdd4},
\url{https://github.com/VladimirReshetnikov/ProveIt}.
The eight reviewed package identities, archive Git blobs and retained
members are recorded in the accompanying source ledgers.
\item Z.-W. Sun, \emph{Further results on Hilbert's Tenth Problem},
Theorems4.12--4.13,
\url{https://arxiv.org/html/1704.03504v7}.
\item Z.-W. Sun, \emph{Ten unknowns for Hilbert's tenth problem over the
integers}, modified Lucas certificate and interpolation,
Theorems3.1 and4.1 and Propositions6.1--6.2,
\url{https://arxiv.org/html/2609.01594v1}.
The exact hypotheses used by the new slice argument are stated in
Section\,\ref{sec:free-lucas-coordinate}; the audit provides a separate proof file.
\item J. P. Jones, D. Sato, H. Wada and D. Wiens,
\emph{Diophantine representation of the set of prime numbers},
American Mathematical Monthly83 (1976),449--464, especially Definition3.7,
Lemma3.8 and Section3.9.
\url{https://www.cs.umd.edu/~gasarch/BLOGPAPERS/Jonesh10.pdf}.
The complete copyrighted article is not redistributed in this archive.
\item The inherited C222 numerical-calculus endpoint and constructive
Diophantine compiler API are retained under
\texttt{direct\_audit/inherited/}, with their own earlier exact repository
pin. These source snapshots expose the interfaces; their transitive
imports are repository dependencies, not newly proved content.
\item Full auxiliary mathematical proofs appear in the component
\texttt{PROOF.md}, \texttt{REVIEW.md}, \texttt{AUDIT.md}, and named theorem
files. Read them together with the coefficient and source manifests.
\end{enumerate}
\paragraph{Reuse.}
Preserved upstream notices identify inherited material. Newly authored
research text, checkers and polynomial data are supplied for incorporation
into ProveIt. The packet asserts neither ownership of inherited material
nor permission to redistribute third-party articles beyond their licenses.
